\section{Trabalhos Relacionados}

Proença e Simões \cite{proenca2020taco} introduziram o TACO com 60 categorias pertencentes a 28 supercategorias. Reconhecendo a ambiguidade visual e o desbalanceamento, os autores avaliaram Mask R-CNN em TACO-1, sem classes, e TACO-10, com nove supercategorias frequentes e uma classe residual. O agrupamento de rótulos no TACO, portanto, já foi explorado. O presente estudo distingue-se pela definição de três níveis ordenados, pela inclusão de um nível intermediário por material e pela reavaliação hierárquica das mesmas detecções para separar localização e semântica.

Córdova et al. \cite{cordova2022litter} compararam Faster R-CNN, Mask R-CNN, RetinaNet, EfficientDet e YOLOv5 em TACO e PlastOPol, relatando dificuldades com objetos pequenos e cenários naturais complexos. O estudo evidencia a relevância de detectores YOLO leves, mas os resultados não são diretamente comparáveis aos obtidos neste trabalho, pois diferem em taxonomia, partição, versão dos modelos e protocolo experimental. O objetivo da presente investigação não consiste em estabelecer o estado da arte, mas em mensurar uma variável experimental sob condições controladas.

YOLO surgiu como detector de estágio único orientado a inferência em tempo real \cite{redmon2016yolo}. O YOLO26 unifica modelos atuais da Ultralytics e oferece uma cabeça \textit{end-to-end} sem NMS, além de mecanismos de atribuição voltados a objetos pequenos \cite{jocher2026yolo26}. A variante nano foi selecionada por sua compatibilidade com uma GPU de 6~GB; não foi realizada comparação entre arquiteturas.

Em aprendizado visual, \emph{cauda longa} descreve a distribuição de frequência dos rótulos: poucas classes de ``cabeça'' concentram grande parte das amostras, enquanto muitas classes da ``cauda'' têm poucos exemplos \cite{cui2019classbalanced}. O termo se refere, portanto, ao desbalanceamento entre classes, não ao tamanho dos objetos nas imagens. Em detecção e segmentação de grande vocabulário, o LVIS demonstrou que essa escassez está associada a desempenho inferior nas categorias pouco representadas \cite{gupta2019lvis}. Perdas balanceadas por classe constituem uma estratégia de mitigação \cite{cui2019classbalanced}. No presente estudo, amostragem e função de perda foram mantidas inalteradas para evitar fatores experimentais adicionais; o desempenho foi mensurado nos estratos raro, médio e frequente.

Diagnósticos de detectores também requerem a separação entre classificação, localização, duplicatas, fundo e omissões. Essa motivação é formalizada pelo TIDE \cite{bolya2020tide}. A decomposição adotada utiliza categorias análogas, mas emprega um procedimento próprio de casamento guloso e deve ser interpretada como análise descritiva no limiar selecionado.
